\documentclass[11pt]{article}
\usepackage[margin=1in]{geometry}
\usepackage[T1]{fontenc}
\usepackage[utf8]{inputenc}
\usepackage{lmodern,microtype,amsmath,amssymb,amsthm,mathtools}
\usepackage[colorlinks=true,linkcolor=blue,citecolor=blue,urlcolor=blue]{hyperref}
\hypersetup{pdftitle={Polynomial total mistake bounds for mean-field Bayesian tanh networks},pdfauthor={ChatGPT-6 Astra; Kaarel Hanni; Dmitry Vaintrob}}
\newtheorem{theorem}{Theorem}
\newtheorem{lemma}[theorem]{Lemma}
\newtheorem{corollary}[theorem]{Corollary}
\newcommand{\R}{\mathbb R}
\newcommand{\E}{\mathbb E}
\newcommand{\Pp}{\mathbb P}
\newcommand{\KL}{\operatorname{KL}}
\title{Polynomial total mistake bounds for\protect\\mean-field Bayesian tanh networks}
\author{ChatGPT-6 Astra \and Kaarel H\"anni \and Dmitry Vaintrob\thanks{Authors are listed in alphabetical order.}}
\date{6 November 2026}
\begin{document}
\maketitle

\begin{abstract}
We study shallow Bayesian tanh networks with independent Gaussian weights
at mean-field scaling. Training conditions on a signed output margin
greater than $1/2$ at each observed input. The initial prior has no output
restriction at unobserved inputs. If a Boolean target on $d$ bits has a
width-$C$ teacher with bounded weights and margin greater than $1/2$,
then, at every sufficiently large learner width, the total number of
online confidence mistakes is $O(d^4+C^4)$, with the weight bound and
positive prior prefactor fixed. The input order may be arbitrary or
adaptive. A confidence mistake means that the posterior assigns at most
a fixed probability, such as $2/3$, to the correct sign; any fixed test
margin below $1/2$ is also allowed. Fresh posterior sampling satisfies
corresponding expected and high-probability total error bounds.
The proof identifies the large-width posterior with successive
relative-entropy projections and shows that each bad prediction consumes
a fixed amount of a bounded entropy budget. We give an explicit Gaussian
mixture certificate for that budget. The width threshold is not
quantified, and the result is not uniform as the prior prefactor tends to
zero. The absence of an initial global output-margin restriction is
essential to the convex projection argument.
\end{abstract}

\section{Model and result}

A small target can admit a polynomial IID sample bound without admitting
a polynomial bound on total mistakes along an arbitrary input sequence.
Here we prove the stronger online conclusion for a Bayesian neural-network
model that requires a positive output margin on observed examples, but
imposes no margin restriction at unobserved inputs before learning starts.

Fix $d\ge1$, let $X=\{-1,1\}^d$, and write $D=d+1$, $M=2^d$, and
$\bar x=(1,x)\in\R^D$. Fix $\tau>0$. At hidden width $N$, sample
independently
\begin{equation}\label{eq:prior}
 u_j\sim\mathcal N(0,\tau^2 I_D/d),\qquad
 a_j\sim\mathcal N(0,\tau^2),\qquad 1\le j\le N,
\end{equation}
and define
\begin{equation}\label{eq:network}
 F_N(x)=\frac1N\sum_{j=1}^N a_j\tanh(u_j\cdot\bar x).
\end{equation}
The physical readout weights $v_j=a_j/N$ have variance $\tau^2/N^2$.
These are the canonical shallow mean-field variances. Equivalently, in
the normalization with first-layer variance $N\sigma_N^2/d$ and readout
variance $\sigma_N^2/N$, they correspond to $\sigma_N=\tau/\sqrt N$.
We take $N\to\infty$ at fixed $\tau$, not a cold limit at fixed width.

Let $f:X\to\{-1,1\}$ be the target. A history is a subset $S\subseteq X$
whose labels have been revealed. The posterior is precisely the prior
conditioned on
\begin{equation}\label{eq:training}
 f(x)F_N(x)>1/2\qquad(x\in S).
\end{equation}
There is \emph{no initial conditioning} on the outputs elsewhere. For a
test-margin threshold $0\le h<1/2$, define
\[
 p_{N,S,h}(x)=\Pp\{f(x)F_N(x)>h\mid \eqref{eq:training}\}.
\]
Fix $0<\theta<1$. An online confidence mistake means
$p_{N,S,h}(x)\le\theta$. Examples are $h=0$, which tests the sign, and
$h=1/4$, which requires a smaller margin than training. We may take
$\theta=2/3$ in either case.

After a prediction, the true label is revealed and the input joins $S$.
The next input may be chosen adaptively. Repeated observed inputs cause
no confidence or sampling errors: their posterior signed output is
already greater than $1/2$. Let $\rho_0$ denote the one-neuron law in
\eqref{eq:prior}, and use $\KL(P\|Q)$ for relative entropy with natural
logarithms.

\begin{theorem}[Total online mistakes]\label{thm:main}
Suppose a probability law $\rho_T$ on $(u,a)\in\R^D\times\R$ satisfies
\begin{equation}\label{eq:certificate}
 \KL(\rho_T\|\rho_0)\le K_T<\infty,
 \qquad
 f(x)\E_{\rho_T}[a\tanh(u\cdot\bar x)]>1/2\quad(x\in X).
\end{equation}
Set $c_0=(\log2)/2$ and
\begin{equation}\label{eq:Rh}
 R_h=\frac{16\tau^2 K_T(K_T+c_0)}{(1/2-h)^2}.
\end{equation}
For each fixed $0<\theta<1$ and $0\le h<1/2$, there exists $N_0$ such
that every $N\ge N_0$ gives at most $\lfloor R_h\rfloor$ total
confidence mistakes. The same $N_0$ works for every input order,
including adaptive choices of the next input.

If predictions use the sign of a fresh posterior sample at each round,
then for every $\eta>0$ there exists $N_1$ such that every $N\ge N_1$
satisfies
\begin{align}
 \E[\text{total sign errors}]&\le R_0+\eta,\label{eq:expectation}\\
 \Pp\{\text{total sign errors}\le\lfloor R_0\rfloor\}
 &\ge1-\eta.\label{eq:hp}
\end{align}
These statements are uniform over adaptive input orders as well.
\end{theorem}

The width thresholds may depend on the target and all fixed parameters,
and need not be polynomial. The conclusion concerns every sufficiently
large width, rather than every width capable of representing the teacher.
The confidence threshold affects the width threshold but does not affect
the displayed mistake bound.

\begin{corollary}[A bounded teacher gives a polynomial bound]\label{cor:teacher}
Suppose
\[
 F_T(x)=\sum_{i=1}^C v_i\tanh(u_i\cdot\bar x),\qquad
 f(x)F_T(x)>1/2\quad(x\in X),
\]
where $C\ge1$ and all coordinates of $u_i$ and all $v_i$ have absolute
value at most $B>0$. Put
\begin{equation}\label{eq:sq}
 s=\min\left\{\frac{\tau}{\sqrt d},\frac1{8CB\sqrt D}\right\},
 \qquad q=\frac{ds^2}{\tau^2}\in(0,1],
\end{equation}
and
\begin{equation}\label{eq:KT}
 K_T=\frac{dDB^2+4C^2B^2}{2\tau^2}
       +\frac D2(q-1-\log q).
\end{equation}
Then Theorem~\ref{thm:main} applies with this value of $K_T$. At fixed
$B,\tau,h$, the total confidence mistake bound is $O(d^4+C^4)$.
The same polynomial controls the expected and high-probability sampling
bounds for the sign convention.
\end{corollary}

The result is not uniform as $\tau\to0$. Even \eqref{eq:KT} has an
inverse-$\tau^2$ contribution. The theorem does not supply a total
mistake bound for a subsequent zero-temperature limit.

\section{The large-width posterior is an entropy projection}

For one neuron, define its full output vector and its feature vector by
\[
 Z(u,a)=\big(a\tanh(u\cdot\bar x)\big)_{x\in X}\in\R^M,
 \qquad h(u)=\big(\tanh(u\cdot\bar x)\big)_{x\in X}.
\]
The network output vector is the empirical average of $N$ independent
copies of $Z$. We prove its conditioning limit directly in this
finite-dimensional space.

\begin{lemma}[A smooth strictly convex rate function]\label{lem:rate}
Let
\[
 \Lambda(\xi)=\log\E_{\rho_0}e^{\xi\cdot Z},\qquad
 I(z)=\sup_{\xi\in\R^M}\{\xi\cdot z-\Lambda(\xi)\}.
\]
The function $\Lambda$ is finite and smooth everywhere, has positive
definite Hessian everywhere, and satisfies
\begin{equation}\label{eq:Lambda}
 \frac{\tau^2}{2}\xi^TG\xi\le\Lambda(\xi)
 \le\frac{\tau^2M}{2}\|\xi\|_2^2,
 \qquad G=\E_u[h(u)h(u)^T]\succ0.
\end{equation}
For every $z$, the supremum defining $I(z)$ has a unique maximizer
$\xi_z$, with $\nabla\Lambda(\xi_z)=z$. The function $I$ is finite,
smooth, strictly convex, and coercive, with $\nabla I(z)=\xi_z$ and
$I(z)\ge\|z\|_2^2/(2\tau^2M)$.
\end{lemma}

\begin{proof}
Integrating out the independent Gaussian $a$ gives
\[
 \Lambda(\xi)=\log\E_u\exp\left\{
       \frac{\tau^2}{2}(\xi\cdot h(u))^2\right\}.
\]
Since each component of $h(u)$ lies in $[-1,1]$, this is at most
$\tau^2\|\xi\|_1^2/2$. Jensen gives the lower bound in
\eqref{eq:Lambda}. Boundedness also permits differentiation locally to
all orders under the integral.

The feature vectors $h(u)$ span $\R^M$. Indeed, choose a linear score
taking distinct values on the $M$ inputs. A nonzero constant tanh feature
and $M-1$ tanh features with thresholds between consecutive scores and
steepness tending to infinity converge to an invertible matrix of
constant and step-function columns. At some finite steepness that matrix
is already invertible. Since the Gaussian $u$ law has positive density
everywhere and $h$ is continuous, $\E_u(\xi\cdot h(u))^2>0$ for every
nonzero $\xi$. Thus $G\succ0$.

The law of $Z$ has full linear span and is symmetric under $a\mapsto-a$,
so it is not supported on an affine hyperplane. Its strictly positive
exponential tilts have the same property. Their covariance matrices,
which are the Hessians of $\Lambda$, are positive definite. The lower
quadratic bound makes $\xi\cdot z-\Lambda(\xi)$ tend to minus infinity
as $\|\xi\|\to\infty$. Its unique maximizer therefore exists and solves
$\nabla\Lambda(\xi_z)=z$. The inverse function theorem gives
\[
 \nabla I(z)=\xi_z,\qquad
 \nabla^2I(z)=\big(\nabla^2\Lambda(\xi_z)\big)^{-1}\succ0.
\]
Taking the Legendre transform of the upper quadratic bound proves the
asserted lower bound on $I$.
\end{proof}

For a history $S$, let
\begin{equation}\label{eq:CS}
 C_S=\{z\in\R^M:f(x)z_x\ge1/2\text{ for }x\in S\}.
\end{equation}
The teacher certificate gives a strictly feasible point. The continuous
coercive rate $I$ therefore has a unique minimizer $z_S$ on this closed
convex set.

\begin{lemma}[Conditional concentration]\label{lem:conditioning}
For each history $S$, the posterior full output vector converges in
probability to $z_S$ as $N\to\infty$. The posterior complement of any
fixed neighborhood of $z_S$ has probability decaying exponentially in $N$.
\end{lemma}

\begin{proof}
Cram\'er's theorem applies to the empirical averages of $Z$, since their
cumulant is finite everywhere \cite[Theorem~2.2.30]{dembozeitouni}.
Its good rate function is $I$. The infimum of $I$ over the strict version
of $C_S$ equals $I(z_S)$: mix $z_S$ with the strict teacher mean and let
the mixture weight tend to zero. Coercivity and uniqueness give a
positive rate gap on the part of $C_S$ outside any fixed open
neighborhood of $z_S$. The large-deviation lower bound on the open
conditioning event and upper bound on that closed complement prove the
claim. In particular these conditioning events have positive probability
for sufficiently large $N$.
\end{proof}

Define the one-neuron exponential tilt
\begin{equation}\label{eq:tilt}
 \frac{d\rho_z}{d\rho_0}
   =\exp\{\xi_z\cdot Z-\Lambda(\xi_z)\}.
\end{equation}
Then $\E_{\rho_z}Z=z$ and $\KL(\rho_z\|\rho_0)=I(z)$. For any law
$\rho$ with mean $z$ and finite relative entropy, the density identity
gives
\begin{equation}\label{eq:KLidentity}
 \KL(\rho\|\rho_0)=\KL(\rho\|\rho_z)+I(z).
\end{equation}
Finite entropy makes the required moments integrable, as
\eqref{eq:moment} below also proves. Thus $\rho_z$ is the unique
minimum-entropy law of mean $z$, and $\rho_S:=\rho_{z_S}$ is the unique
entropy projection onto the training constraints. This is the usual
relative-entropy, or Sanov, interpretation of the variational problem.
The actual posterior convergence has already been proved directly by
the finite-dimensional Cram\'er argument.

\section{A bounded potential pays for every bad prediction}

The same teacher law is feasible at every target-consistent history, so
\begin{equation}\label{eq:budget}
 0\le I(z_S)\le K_T.
\end{equation}
For the empty history, $z_\varnothing=0$ and $I(z_\varnothing)=0$.

\begin{lemma}[Projection increment]\label{lem:projection}
If $S'$ extends $S$ by one correct training example, then
\begin{equation}\label{eq:increment}
 I(z_{S'})-I(z_S)\ge\KL(\rho_{S'}\|\rho_S).
\end{equation}
\end{lemma}

\begin{proof}
Since $z_{S'}\in C_S$, first-order optimality for the convex minimization
at $z_S$ gives $\xi_{z_S}\cdot(z_{S'}-z_S)\ge0$. The tilt identity gives
\[
 \KL(\rho_{S'}\|\rho_S)
  =I(z_{S'})-I(z_S)-\xi_{z_S}\cdot(z_{S'}-z_S),
\]
which proves the claim.
\end{proof}

For any $\rho$ with $\KL(\rho\|\rho_0)\le K_T$, the entropy inequality
and $\E_{\rho_0}e^{a^2/(4\tau^2)}=\sqrt2$ imply
\begin{equation}\label{eq:moment}
 \E_\rho a^2\le4\tau^2(K_T+c_0),\qquad c_0=\tfrac12\log2.
\end{equation}
Applying the entropy inequality to truncations of $a^2$ and passing to
the limit also proves moment finiteness. The observable
$H_x=a\tanh(u\cdot\bar x)$ satisfies $H_x^2\le a^2$.

For probability measures $P,Q$ and an observable $H$ with finite second
moments, Cauchy--Schwarz gives
\begin{align}
 |\E_PH-\E_QH|^2
 &\le2(\E_PH^2+\E_QH^2)
       \int(\sqrt{dP}-\sqrt{dQ})^2\notag\\
 &\le2(\E_PH^2+\E_QH^2)\KL(P\|Q).
 \label{eq:hellinger}
\end{align}
The first line factors the density difference into the sum and difference
of square roots. The second uses squared Hellinger distance at most
relative entropy, with squared Hellinger distance normalized as the
displayed integral.

If the next query $x$ has limiting signed output $f(x)(z_S)_x\le h$,
the new training constraint forces $f(x)(z_{S'})_x\ge1/2$. Taking
$P=\rho_{S'}$, $Q=\rho_S$, and $H=f(x)H_x$ in
\eqref{eq:hellinger}, and using \eqref{eq:budget}--\eqref{eq:moment},
gives
\begin{equation}\label{eq:payment}
 \KL(\rho_{S'}\|\rho_S)
 \ge\frac{(1/2-h)^2}{16\tau^2(K_T+c_0)}.
\end{equation}
Each such round increases the potential by this amount. Its total budget
is $K_T$, so at most $R_h$ rounds have $f(x)(z_S)_x\le h$, along every
input order. In particular,
\[
 R_0=64\tau^2K_T(K_T+c_0),\qquad
 R_{1/4}=256\tau^2K_T(K_T+c_0).
\]

\section{From the limiting predictor to the actual posterior}

There are finitely many target-consistent histories: after removing
repeats, each is a subset of the finite domain. For every pair $(S,x)$
with $f(x)(z_S)_x>h$, conditional concentration gives
$p_{N,S,h}(x)\to1$. A common sufficiently large width makes all these
probabilities greater than $\theta$. Every actual confidence mistake
must then be among the limiting bad rounds counted by \eqref{eq:Rh}.
This proves the first part of Theorem~\ref{thm:main}, uniformly over
input orders.

For posterior sampling take $h=0$ and choose the common width so each
limiting-good pair has wrong-sign probability at most $\eta/M$. At bad
pairs use the bound one. Every order has at most $\lfloor R_0\rfloor$
bad rounds and at most $M$ other distinct rounds. Linearity of expectation
gives \eqref{eq:expectation}. A union bound shows that with probability
at least $1-\eta$ no good round makes an error, giving \eqref{eq:hp}.
For adaptive input choices the bounds apply conditionally on the past,
because the next posterior draw is fresh. Repeated observed inputs have
zero error probability. This completes the proof of the theorem.

\section{An explicit polynomial entropy certificate}

To prove Corollary~\ref{cor:teacher}, let $s,q$ be as in \eqref{eq:sq}.
For each teacher neuron $i$, take the independent Gaussian component
\[
 u\sim\mathcal N(u_i,s^2I_D),\qquad
 a\sim\mathcal N(2Cv_i,\tau^2),
\]
and let $\rho_T$ be their uniform mixture. Its mean output is
\[
 \E_{\rho_T}H_x
   =2\sum_i v_i\E_\zeta\tanh((u_i+\zeta)\cdot\bar x),
 \qquad\zeta\sim\mathcal N(0,s^2I_D).
\]
Since tanh is $1$-Lipschitz and $\|\bar x\|=\sqrt D$,
\[
 |\E_{\rho_T}H_x-2F_T(x)|
 \le2\sum_i|v_i|\E|\zeta\cdot\bar x|
 \le2CBs\sqrt D\le1/4.
\]
The signed mean is therefore strictly greater than $3/4$ on the entire
cube. The Gaussian entropy formula bounds each component's entropy by
\[
 \frac{d\|u_i\|^2+(2Cv_i)^2}{2\tau^2}
   +\frac D2(q-1-\log q)
 \le\frac{dDB^2+4C^2B^2}{2\tau^2}
   +\frac D2(q-1-\log q).
\]
Convexity of relative entropy under mixing gives the same bound for
$\rho_T$. No $\log C$ mixture penalty is needed.

At fixed $B,\tau$, the first term is $O(d^2+C^2)$, while
\[
 q-1-\log q\le-\log q
 =\max\left\{0,\log\frac{64C^2B^2D\tau^2}{d}\right\}
 =O(1+\log C).
\]
Because $d(1+\log C)=O(d^2+C^2)$, the certificate satisfies
$K_T=O(d^2+C^2)$ and $R_h=O(d^4+C^4)$ at fixed $B,\tau,h$.
All constants are displayed in \eqref{eq:Rh}--\eqref{eq:KT}.

\section{Scope and the role of the support restriction}

Training defines nested convex constraints on the full mean output
vector. A prediction below the test margin forces a nonzero movement
before the next training margin can be met. Relative entropy pays for
that movement from a bounded budget. This property does not require an
IID input distribution.

If the initial prior is instead restricted to networks with
$|F_N(x)|>1/2$ at every input, the limiting feasible set becomes a union
of margin orthants. It is generally nonconvex: averaging two valid
outputs can create zeros at unobserved inputs. The first-order inequality
in Lemma~\ref{lem:projection} need not hold in that model. We make no
claim here about total mistakes for that globally restricted learner.

Two other distinctions matter. Training uses a fixed positive margin.
For sign-only training at zero margin, the origin lies in the closure of
every feasible set, so this rate minimizer no longer encodes the labels.
Also, the confidence guarantee tests a margin strictly smaller than the
training margin. The difference $1/2-h$ pays for each error in
\eqref{eq:payment}.

The abstract proof extends to bounded activations whose ridge feature
vectors span the finite input set, provided a finite-entropy strict
teacher law exists. A Lipschitz bound also permits the Gaussian mixture
certificate, with modified constants. This does not establish the same
theorem for ReLU: an unbounded activation makes a neuron output a product
of unbounded Gaussian quantities, and the global cumulant and
entropy-to-second-moment estimates require a different argument.

\begin{thebibliography}{9}
\bibitem{dembozeitouni}
Amir Dembo and Ofer Zeitouni.
\emph{Large Deviations Techniques and Applications}, second edition.
Springer, 1998.
\end{thebibliography}
\end{document}
